\documentclass[11pt]{article}
\usepackage[a4paper,margin=1in]{geometry}
\usepackage{amsmath,amssymb,amsthm}
\usepackage{hyperref}
\usepackage{parskip}
\usepackage{enumitem}

\newtheorem{theorem}{Theorem}
\newtheorem{proposition}[theorem]{Proposition}

\title{Recursive Language Model Children as Options:\\
       a Sutton-Precup-Singh correspondence}
\author{Pablo Leyva}
\date{May 15, 2026}

\begin{document}
\maketitle

\section*{Setup: the Sutton-Precup-Singh options framework}

Recall the options framework introduced by Sutton, Precup, and Singh (1999,
\emph{Artificial Intelligence} 112). An MDP is a tuple
$(\mathcal{S}, \mathcal{A}, P, R, \gamma)$. An \emph{option} is a triple
$o = (I_o, \pi_o, \beta_o)$ where:
\begin{itemize}
\item $I_o \subseteq \mathcal{S}$ is the \emph{initiation set} --- states from
      which $o$ can be invoked,
\item $\pi_o : \mathcal{S} \to \Delta(\mathcal{A})$ is the option's
      \emph{internal policy},
\item $\beta_o : \mathcal{S} \to [0, 1]$ is the \emph{termination function} ---
      the probability that $o$ terminates upon reaching state $s$.
\end{itemize}
A high-level controller selects options in $I_o \subseteq \mathcal{S}$;
each option runs to termination, after which the controller selects another
option. The induced process is a \emph{semi-Markov decision process} (SMDP).

For the SMDP, define the option-value
\[
Q_\Omega(s, o) = \mathbb{E}\Bigl[\sum_{t=0}^{K-1} \gamma^t r_{t+1} + \gamma^K
                 \max_{o' \in \mathcal{O}_{s_K}} Q_\Omega(s_K, o')
                 \,\Big|\, s_0 = s,\ o_0 = o\Bigr],
\]
where $K$ is the (random) duration of $o$, $r_{t+1}$ are intra-option rewards,
and $\mathcal{O}_{s_K}$ is the set of options available at termination state
$s_K$. SPS Theorem~1 (1999, Theorem 1) shows that SMDP Q-learning with
options converges with probability 1 to $Q_\Omega^*$ under standard
Robbins-Monro stepsize and ergodicity conditions.

\section*{Mapping: \texttt{rlm\_query} calls as options}

Recall the RLM setup of \texttt{02-math-deep-dive.md}: the agent operates in a
Python REPL with primitives \texttt{FINAL}, \texttt{FINAL\_VAR},
\texttt{rlm\_query}. Each \texttt{rlm\_query(p)} call spawns a child rollout
under the same policy $\pi_\theta$, returning the child's \texttt{FINAL}
answer to the parent's REPL.

We propose the following correspondence. Let $\mathcal{S}_{\mathrm{REPL}}$ be
the set of REPL states (Python heap snapshots plus prompt history). Let
$\mathcal{P}$ be the (possibly infinite) set of prompts that can appear as the
argument of a \texttt{rlm\_query} call. For each prompt $p \in \mathcal{P}$,
define option $o_p = (I_p, \pi_{o_p}, \beta_p)$ by:
\begin{itemize}
\item $I_p = \{s \in \mathcal{S}_{\mathrm{REPL}} : \pi_\theta\text{ may emit } \texttt{rlm\_query}(p) \text{ at } s\}$,
\item $\pi_{o_p}(\cdot \mid s) = \pi_\theta(\cdot \mid s, \text{prompt} = p)$ ---
      the same shared policy, conditioned on the child prompt $p$,
\item $\beta_p(s) = \mathbf{1}[s \text{ contains a \texttt{FINAL}(.)} \text{ invocation}]$.
\end{itemize}

Primitive actions (one-token emissions that are not \texttt{rlm\_query} calls)
are kept as-is: each token-emission action is a degenerate option with
$\beta = 1$ everywhere and duration $K = 1$. This makes the entire RLM
session an SMDP over the high-level alphabet
$\mathcal{A}_{\Omega} = \mathcal{A}_{\mathrm{primitive}} \cup \{o_p : p \in \mathcal{P}\}$.

The intra-option discounted reward of $o_p$ is the discounted sum of (zero)
intra-child rewards plus the discounted child \emph{advantage} accumulated
through advantage inheritance. In the undiscounted ($\gamma = 1$) limit,
$Q_\Omega(s, o_p)$ reduces to the expected root reward $r_g$ given that the
agent invoked $o_p$ from state $s$ and then continued under $\pi_\theta$.

\section*{Theorem: convergence carry-over}

\begin{theorem}[SMDP Q-learning convergence for RLM SMDPs]
\label{thm:smdp}
Let $\theta$ parameterise an autoregressive policy $\pi_\theta$ acting in the
RLM SMDP defined by the mapping above. Let $\widehat Q_\Omega$ be the
SMDP-Q-learning estimator,
\[
\widehat Q_\Omega(s, o) \leftarrow \widehat Q_\Omega(s, o) + \alpha_n
   \Bigl[ R_o + \gamma^{K_o} \max_{o' \in \mathcal{O}_{s'}} \widehat Q_\Omega(s', o')
   - \widehat Q_\Omega(s, o)\Bigr],
\]
with the conditions:
\begin{enumerate}[label=(\roman*)]
\item Robbins-Monro stepsizes: $\sum_n \alpha_n = \infty$, $\sum_n \alpha_n^2 < \infty$.
\item Every $(s, o)$ pair is visited infinitely often.
\item Bounded reward, finite SMDP option set at each state, finite expected
      option duration.
\end{enumerate}
Then $\widehat Q_\Omega \to Q^*_\Omega$ with probability 1.
\end{theorem}

\begin{proof}[Proof sketch]
The RLM SMDP, under the mapping of the previous section, satisfies all
hypotheses of Sutton-Precup-Singh 1999, Theorem~1: it is an SMDP with finite
option set at each state (the vocabulary is finite, so $\mathcal{P}$ is
countable, and only those tokens producible by $\pi_\theta$ at $s$ are
available); option durations are finite in expectation (the harness imposes
a max number of turns per child rollout); rewards are bounded ($r_g \in [0, 1]$
for rubric scores). Conditions (i)-(iii) are precisely the SPS hypotheses.
The conclusion follows by direct application of SPS Theorem~1, with no new
argument required --- the mapping is a faithful translation, and the SPS
proof goes through unchanged.
\end{proof}

\begin{proposition}[Hierarchical Bellman equation for RLM]
For the optimal SMDP value $V^*_\Omega(s) = \max_o Q^*_\Omega(s, o)$, we have
the hierarchical Bellman equation
\[
V^*_\Omega(s) = \max_{o \in \mathcal{O}_s}
   \mathbb{E}\bigl[R_o + \gamma^{K_o} V^*_\Omega(s')\bigr].
\]
This recovers the recursive subtree loss of concept~11 as the policy-gradient
analogue: the recursion in the loss is exactly the recursion in the
hierarchical Bellman backup.
\end{proposition}

\begin{proof}[Proof]
Standard Bellman optimality applied to the SMDP. The connection to the
recursive subtree loss is by inspection: both expressions are recursive over
the tree of \texttt{rlm\_query} calls with the same accumulating structure.
\end{proof}

\section*{Discussion}

\paragraph{What is gained from the connection.}
\begin{itemize}
\item \textbf{Free convergence guarantee.} Theorem~\ref{thm:smdp} promotes
      RLM-Q-learning from heuristic to convergent under standard conditions ---
      no new proof required.
\item \textbf{Hierarchical Bellman equation.} The recursive subtree loss of
      concept~11 is no longer just a notational convenience; it is the
      policy-gradient counterpart of a value-function recursion that has been
      studied for over two decades.
\item \textbf{Options as actions for higher-level planning.} Once
      \texttt{rlm\_query} calls are options, one can compose them into
      \emph{multi-level} options (an option that invokes a sequence of
      \texttt{rlm\_query} calls), giving a clean path to deeper hierarchical
      planning without changing the training algorithm.
\item \textbf{Intra-option learning.} SPS 1999 introduced intra-option
      Q-learning, which updates Q-values from primitive transitions inside an
      option (not just at termination). This translates directly to RLM:
      every token within a child rollout could supply a Q-update for the
      enclosing option, accelerating learning.
\item \textbf{Interruptibility.} SPS 1999 also introduced interruption: stop
      an option early if a higher-value option becomes available. For RLM, this
      is a principled way to abort a child rollout that is going badly without
      having to wait for its \texttt{FINAL}.
\end{itemize}

\paragraph{What does \emph{not} immediately carry over.} The shared-policy
constraint of the paper means the \emph{internal policies} of all options are
the \emph{same} $\pi_\theta$ (modulo prompt conditioning). The SPS framework
allows arbitrarily different internal policies. This shared-policy
restriction does not break SMDP convergence (Theorem~\ref{thm:smdp} still
holds) but it does mean some SPS extensions --- e.g., learning option
\emph{policies} independently from the high-level controller --- collapse to
trivial cases.

\paragraph{Recommended next step.} Formalise intra-option Q-learning in the
RLM setting and check empirically whether it converges faster than the
inheritance-only estimator of the blog. This is the most actionable corollary.

\end{document}
